% Part:sets-functions-relations
% Chapter: size-of-sets
% Section: comparing-sizes

\documentclass[../../../include/open-logic-section]{subfiles}

\begin{document}

\olfileid[fa-IR]{sfr}{siz}{car}

\olsection{مجموعه‌هایی با اندازه‌های متفاوت و قضیهٔ کانتور}

\begin{explain}
بیان دقیقی از این اندیشه به دست دادیم که دو مجموعه اندازه‌ای یکسان
دارند. همچنین می‌توانیم این اندیشه را که یک مجموعه از دیگری کوچک‌تر
است به‌دقت بیان کنیم. تعریف ما از «کوچک‌تر از (یا هم‌توان با)» به‌جای
!!a{bijection} میان مجموعه‌ها، به !!a{injection} از مجموعهٔ نخست به
مجموعهٔ دوم نیاز دارد. اگر چنین تابعی وجود داشته باشد، اندازهٔ مجموعهٔ
نخست کوچک‌تر یا مساوی اندازهٔ مجموعهٔ دوم است. به‌طور شهودی، وجود
!!a{injection} از یک مجموعه به مجموعه‌ای دیگر تضمین می‌کند که تعدادِ
!!{element}s در برد تابع دست‌کم برابر با تعداد آن‌ها در دامنه است، زیرا
هیچ دو مورد از !!{element}s واقع در دامنه به یک !!{element} در برد
نگاشته نمی‌شوند.
\end{explain}

\begin{defn}
مجموعهٔ $A$ از~$B$ \emph{بزرگ‌تر نیست}، که آن را با $\cardle{A}{B}$
می‌نویسیم، اگر و تنها اگر !!a{injection} به‌صورت $f \colon A \to B$
وجود داشته باشد.
\end{defn}

آشکار است که این رابطه بازتابی و تراگذری است، اما متقارن نیست (اثبات
این مطلب به‌عنوان تمرین واگذار می‌شود). همچنین می‌توانیم مفهومی معرفی
کنیم که می‌گوید یک مجموعه به‌طور اکید از دیگری کوچک‌تر است.

\begin{defn}
مجموعهٔ $A$ از~$B$ \emph{کوچک‌تر است}، که آن را با $\cardless{A}{B}$
می‌نویسیم، اگر و تنها اگر !!a{injection}~$f\colon A \to B$ وجود داشته
باشد، اما هیچ !!{bijection}~$g\colon A
\to B$ وجود نداشته باشد؛ یعنی $\cardle{A}{B}$ و $\cardneq{A}{B}$.
\end{defn}

آشکار است که این رابطه غیربازتابی و تراگذری است. (اثبات این مطلب
به‌عنوان تمرین واگذار می‌شود.) با این نمادگذاری می‌توانیم بگوییم
مجموعهٔ $A$ !!{enumerable} است اگر و تنها اگر $\cardle{A}{\Nat}$، و
$A$ !!{nonenumerable} است اگر و تنها اگر $\cardless{\Nat}{A}$. این
نمادگذاری به ما اجازه می‌دهد
\oliflabeldef{sfr:siz:nen-alt:thm:nonenum-pownat}{%
\olref[sfr][siz][nen-alt]{thm:nonenum-pownat}
را به‌صورت مشاهدهٔ
$\cardless{\Nat}{\Pow{\Nat}}$ بازگو کنیم}{%
\olref[sfr][siz][nen]{thm:nonenum-pownat}
را به‌صورت مشاهدهٔ $\cardless{\PosInt}{\Pow{\PosInt}}$ بازگو کنیم}. در واقع،
\citet{Cantor1892} ثابت کرد که این نکته \emph{کاملاً کلی} است:

\begin{thm}[کانتور]\ollabel{thm:cantor}
$\cardless{A}{\Pow{A}}$ برای هر مجموعهٔ $A$ برقرار است.
\end{thm}

\begin{proof}
نگاشت $f(x) = \{x\}$، !!a{injection} به‌صورت $f \colon A \to \Pow{A}$
است، زیرا اگر $x \neq y$، آنگاه بنا بر اصل امتداد،
$\{x\} \neq \{y\}$ نیز برقرار است و بنابراین $f(x) \neq f(y)$.
پس $\cardle{A}{\Pow{A}}$.

\begin{editorial}
اگر \olref[nen]{sec} موجود باشد برهان تفصیلی را می‌آوریم، وگرنه برهان
سریع‌تری را که با \olref[nen-alt]{sec} مطابقت دارد ارائه می‌کنیم.
\end{editorial}

\oliflabeldef{sfr:siz:nen:sec}{%
  اکنون نشان می‌دهیم که هیچ تابع !!a{surjective} به‌صورت~$g\colon A \to
  \Pow{A}$ نمی‌تواند وجود داشته باشد، چه رسد به تابعی !!a{bijective}؛
  و ازاین‌رو $\cardneq{A}{\Pow{A}}$. فرض کنید $g\colon A \to \Pow{A}$.
  چون $g$ تام است، هر $x \in A$ به زیرمجموعه‌ای مانند $g(x)
  \subseteq A$ نگاشته می‌شود. می‌توانیم نشان دهیم که $g$ نمی‌تواند
  پوشا باشد. برای این کار، زیرمجموعه‌ای مانند~$\overline{A} \subseteq A$
  تعریف می‌کنیم که بنا بر تعریف نمی‌تواند در بردِ~$g$ باشد. بگذارید
  \[
  \overline{A} = \Setabs{x \in A}{x \notin g(x)}.
  \]
  چون $g(x)$ برای همهٔ $x \in A$ تعریف شده است، $\overline{A}$ به‌روشنی
  زیرمجموعه‌ای خوش‌تعریف از~$A$ است. اما نمی‌تواند در بردِ~$g$ باشد.
  $x \in A$ را دلخواه بگیرید؛ نشان می‌دهیم $\overline{A} \neq
  g(x)$. اگر $x \in g(x)$، آنگاه شرط $x \notin g(x)$ را برآورده
  نمی‌کند؛ پس بنا بر تعریفِ~$\overline{A}$، داریم $x \notin
  \overline{A}$. اگر $x \in \overline{A}$، باید ویژگی معرفِ
  $\overline{A}$ را برآورده کند؛ یعنی $x \in A$ و $x \notin
  g(x)$. چون $x$ دلخواه بود، این نشان می‌دهد که برای هر $x \in
  \overline{A}$، $x \in g(x)$ اگر و تنها اگر $x \notin \overline{A}$؛
  پس $g(x) \neq \overline{A}$. به‌عبارت دیگر، $\overline{A}$ نمی‌تواند
  در بردِ~$g$ باشد و این با فرض پوشابودنِ~$g$ تناقض دارد.}{هنوز باید
  نشان دهیم $\cardneq{A}{\Pow{A}}$. به قصد خلف، فرض کنید
  $\cardeq{A}{\Pow{A}}$؛ یعنی یک !!{bijection} به‌صورت
  $g \colon A \to \Pow{A}$ وجود دارد. اکنون در نظر بگیرید:
  \[
    D = \Setabs{x \in A}{x \notin g(x)}
  \]
  توجه کنید که $D \subseteq A$؛ پس $D \in \Pow{A}$. چون $g$
  !!a{bijection} است، عضوی مانند $y \in A$ وجود دارد که $g(y) = D$. اما اکنون
  داریم:
  \[
    y \in g(y) \text{ اگر و تنها اگر } y \in D \text{ اگر و تنها اگر } y \notin g(y).
  \]
  این تناقض است؛ پس $\cardneq{A}{\Pow{A}}$.}{}
\end{proof}

\begin{explain}
\oliflabeldef{sfr:siz:nen:thm:nonenum-pownat}{مقایسهٔ برهان
  \olref{thm:cantor} با برهانِ
  \olref[nen]{thm:nonenum-pownat} آموزنده است. آنجا نشان دادیم که برای
  هر فهرستِ $Z_1$، $Z_2$، \dots از زیرمجموعه‌های~$\PosInt$ می‌توان
  مجموعه‌ای از اعداد به‌صورت~$\overline{Z}$ ساخت که تضمین شده است در
  فهرست نباشد. این مجموعه از آن رو در فهرست نبود که برای هر $n \in
  \PosInt$، $n \in Z_n$ اگر و تنها اگر $n \notin \overline{Z}$. به
  این ترتیب، همیشه عددی وجود دارد که !!a{element} یکی از $Z_n$ یا
  $\overline{Z}$ است، اما عضو آن دیگری نیست. در اینجا نیز از همان اندیشه
  پیروی می‌کنیم، جز اینکه شاخص‌های~$n$ اکنون از جملهٔ !!{element}s در~$A$ هستند،
  نه عضوهای~$\PosInt$. مجموعهٔ $\overline{A}$ چنان تعریف می‌شود که
  با~$g(x)$ برای هر $x \in A$ متفاوت باشد، زیرا $x \in g(x)$ اگر و تنها
  اگر $x \notin \overline{A}$. باز همواره !!a{element}ی از~$A$ وجود دارد
  که !!a{element} یکی از $g(x)$ و $\overline{A}$ است، اما عضو آن دیگری
  نیست. همان‌گونه که $\overline{Z}$ ازاین‌رو نمی‌تواند در فهرستِ
  $Z_1$، $Z_2$، \dots باشد، $\overline{A}$ نیز نمی‌تواند در بردِ~$g$
  باشد.}{}

\oliflabeldef{sfr:siz:nen-alt:thm:nonenum-pownat}{مقایسهٔ برهان
  \olref{thm:cantor} با برهانِ
  \olref[nen-alt]{thm:nonenum-pownat} آموزنده است. آنجا نشان دادیم که
  برای هر فهرستِ $N_0$، $N_1$، $N_2$، \dots از زیرمجموعه‌های~$\Nat$
  می‌توان مجموعه‌ای از اعداد به‌صورت~$D$ ساخت که تضمین شده است در
  فهرست نباشد. این مجموعه از آن رو در فهرست نبود که
  $n \in N_n$ اگر و تنها اگر $n \notin
  D$، و این برای هر $n \in \Nat$ برقرار است. در اینجا نیز از همان اندیشه پیروی می‌کنیم،
  جز اینکه شاخص‌های~$n$ اکنون از جملهٔ !!{element}s در~$A$ هستند، نه عضوهای~$\Nat$.
  مجموعهٔ $D$ چنان تعریف می‌شود که با~$g(x)$ برای هر $x
  \in A$ متفاوت باشد، زیرا $x \in g(x)$ اگر و تنها اگر $x \notin D$.}{}

این برهان را همچنین می‌توان با برهان پارادوکس راسل مقایسه کرد:
\olref[sfr][set][rus]{thm:russells-paradox}. در واقع قضیهٔ کانتور الهام‌بخش
پارادوکس خود راسل بود.
\end{explain}

\begin{prob}
  نشان دهید هیچ !!a{injection} به‌صورت $g\colon \Pow{A} \to
  A$ برای هیچ مجموعهٔ~$A$ نمی‌تواند وجود داشته باشد. راهنمایی: فرض کنید
  $g\colon \Pow{A} \to A$
  !!{injective} است. $D = \Setabs{g(B)}{B \subseteq A \text{ و
  } g(B) \notin B}$ را در نظر بگیرید. بگذارید $x = g(D)$. با استفاده
  از این واقعیت که $g$ !!{injective} است، تناقضی به دست آورید.
\end{prob}

\end{document}
